% Estilo compartilhado pelas versoes PT e EN
\usepackage[a4paper,top=1.3cm,bottom=1.3cm,left=1.5cm,right=1.5cm]{geometry}
\usepackage[T1]{fontenc}
\usepackage[utf8]{inputenc}
\usepackage{lmodern}
\usepackage{titlesec}
\usepackage{enumitem}
\usepackage{xcolor}
\usepackage[hidelinks]{hyperref}
\usepackage{tabularx}

\definecolor{linkblue}{RGB}{20,70,150}
\newcommand{\link}[2]{\href{#1}{\textcolor{linkblue}{#2}}}

\pagestyle{empty}
\setlength{\parindent}{0pt}
\linespread{1.04}

% Titulos de secao em caixa alta com linha abaixo
\titleformat{\section}{\large\bfseries\raggedright}{}{0em}{\MakeUppercase}[\vspace{1pt}\titlerule]
\titlespacing*{\section}{0pt}{10pt}{4pt}

% \entrada{titulo em negrito}{detalhe}{data a direita}
\newcommand{\entrada}[3]{%
  \par\vspace{4pt}\noindent\textbf{#1}#2\hfill#3\par\vspace{1pt}}

% Itens com ponto medio, como no modelo
\newenvironment{itens}{\begin{itemize}[leftmargin=0pt,label={$\cdot$},itemsep=1pt,topsep=1pt,parsep=0pt,itemindent=8pt,labelsep=4pt]}{\end{itemize}}

% Linha "Rotulo   conteudo" para competencias
\newcommand{\linha}[2]{\noindent\begin{tabularx}{\linewidth}{@{}p{4.0cm}@{}X@{}}\textbf{#1} & #2\end{tabularx}\par\vspace{0pt}}
